% ===================================================================
%  TAULUKKO N:o 3 — Lapsuus ja nuoruus.
%  Traduction 2026 d'apres texto/io, controlee sur texto/fr et
%  texto/sv. Consigne : voir texto/fi/10-taulukko-01.tex.
%
%  UNE INVERSION EVITEE D'AVANCE, au bloc c1-04-1. L'ido ecrit
%  « circum la spiro (7) di la kloshturmo (5) », le possede d'abord ;
%  le genitif finnois preposerait — « kellotornin (5) huippu (7) » —
%  et les deux renvois changeraient de rang. On ecrit donc
%  « huipun (7) ymparilla, joka on kellotornissa (5) » : une relative
%  a la place du genitif, et l'ordre est rendu.
% ===================================================================

%%K t03-apar-1 apar
\VUsaut{3.0mm}
\VUtitre{15.0pt}{60}{TAULUKKO N:o 3}
\VUsaut{1.6mm}
\VUtitre{11.4pt}{0}{Lapsuus ja nuoruus.}
\VUsaut{3.4mm}

%%K t03-apar-2 apar
(3\textsuperscript{s} ja 4\textsuperscript{s} taulukko ovat niin laaditut, että
ne neljässä \VUgras{perhekohtauksessa} esittävät olennaiset käsitteet
\VUgras{säästä}, \VUgras{iästä}, \VUgras{perheestä}, \VUgras{pukeutumisesta} ja
\VUgras{säädystä}.)
\VUsaut{4.0mm}

%%K t03-tit-1 sub
\VUcentre{12.0pt}{0}{\textit{Ensimmäinen kohtaus.}}
\VUsaut{3.0mm}

%%K t03-tit-2 sub
\VUpk{10.2pt}{40}{Kastejuhla kylässä.}
\VUsaut{2.4mm}

%%K t03-tit-3 sub
\VUpk{10.2pt}{40}{Kevät.}
\VUsaut{3.0mm}

%%K t03-c1-01-1 p
1. --- Olemme \VUgras{keväässä}, \VUgras{maaliskuun}, \VUgras{huhtikuun} tai
\VUgras{toukokuun} kuukausina, \VUgras{viikon} päivinä, \VUgras{maanantaina},
\VUgras{tiistaina} tai \VUgras{keskiviikkona}. Kello on \VUgras{kymmenen}
aamulla.

%%K t03-c1-02-1 p
2. --- \VUgras{Sää} on kaunis, \VUgras{ilma} puhdas. Aurinko paistaa. Muutamia
\VUgras{pieniä pilviä} \textsuperscript{(2)} nousee siniselle \VUgras{taivaalle}
\textsuperscript{(1)}.

%%K t03-c1-03-1 p
3. --- \VUgras{Puissa} on tuskin lainkaan \VUgras{lehtiä}. Hoikissa
\VUgras{poppeleissa} \textsuperscript{(3)} ja korkeassa \VUgras{plataanissa}
\textsuperscript{(17)} on tähän asti vain silmuja.

%%K t03-c1-04-1 p
4. --- \VUgras{Pääskyset} \textsuperscript{(4)} palaavat ja lentävät iloisesti
visertäen \VUgras{huipun} \textsuperscript{(7)} ympärillä, joka on
\VUgras{kellotornissa} \textsuperscript{(5)}, joka kruunaa kylän
\VUgras{kirkkoa} \textsuperscript{(6)}.

%%K t03-tit-4 sub
\VUsaut{3.4mm}
\VUpk{10.2pt}{40}{Kirkko.}
\VUsaut{3.0mm}

%%K t03-c1-05-1 p
5. --- Hyvin yksinkertainen on tuo kirkko suurine \VUgras{pääportaaleineen}
\textsuperscript{(13)} ja paksuine \VUgras{kelloineen} \textsuperscript{(8)}.
Pieni \VUgras{viisari} \textsuperscript{(9)} on numerolla X, se näyttää
\VUgras{tunnit}. \VUgras{Suuri} \textsuperscript{(10)} on numerolla XII
(keskipäivä); se näyttää \VUgras{minuutit}.

%%K t03-c1-06-1 p
6. --- Kellotornissa soivat \VUgras{kirkonkellot} \textsuperscript{(11)}
iloisesti. Katon huipulla on pieni kivinen \VUgras{risti}
\textsuperscript{(12)}.

%%K t03-c1-07-1 p
7. --- Kirkossa ei ole vain suurta \VUgras{pääportaalia}
\textsuperscript{(13)}, siinä on myös pieni sivuovi. Pienestä ovesta herra
\VUgras{kirkkoherra} \textsuperscript{(15)} \VUgras{kauhtanaansa}
pukeutuneena tulee ulos \VUgras{sakastista} \textsuperscript{(14)} ja menee
\VUgras{pappilaan} \textsuperscript{(19)}.

%%K t03-c1-08-1 p
8. --- Hänen vieressään, tässä on \VUgras{maidonmyyjä} \textsuperscript{(24)}
\VUgras{aasinsa} \textsuperscript{(23)} kanssa. Hän palaa kaupungista, jonne hän
vei hyvää maitoa lapsille.

%%K t03-tit-5 sub
\VUsaut{4.0mm}
\VUpk{10.2pt}{40}{Hedelmätarha.}
\VUsaut{3.4mm}

%%K t03-c1-09-1 p
9. --- Herra Aliboron (aasi) on aivan tyytyväinen tuohon aamuretkeen. Hän
hengittää ihastuksissaan ilmaa, jonka on balsamoinut niiden \VUgras{kukkien}
\textsuperscript{(20)} tuoksu, jotka peittävät läheisen \VUgras{hedelmätarhan}
\VUgras{hedelmäpuita} \textsuperscript{(18)}. Aivan ruusunpunaisia ja valkoisia
ovat \VUgras{persikkapuut} \textsuperscript{(18)} ja \VUgras{aprikoosipuut}
\textsuperscript{(19)}, ja ne antavat toivoa hyvästä sadosta.

%%K t03-c1-10-1 p
10. --- Sillä aikaa pienet \VUgras{mehiläiset} \textsuperscript{(22)} menevät
sinne \VUgras{mehiläispesästä} \textsuperscript{(21)} suristen ryöstämään
suloista \VUgras{hunajaansa}.

%%K t03-c1-11-1 p
11. --- Mutta mitä tapahtuu? Tässä tulevat kylän lapset juosten ja karkottavat
säikähtäneet \VUgras{hanhet} \textsuperscript{(25)}. Se on saattueen lähtö
kirkosta notaarin pojan \VUgras{kasteen} jälkeen.

%%K t03-c1-12-1 p
12. --- Pikku Jaakob herää \VUgras{kehdossaan} \textsuperscript{(26)}
kirkonkellojen iloiseen soittoon, ja hänen sisarensa Magdaleena, joka myös
haluaa nähdä kastesaattueen, pyytää äitiään tulemaan kampaamaan hänen hiuksensa
ja pukemaan hänet talon kynnyksellä.

%%K t03-c1-13-1 p
13. --- Äiti suostuu. Hän panee päälleen \VUgras{päähuivinsa}
\textsuperscript{(28)}, \VUgras{liivinsä} \textsuperscript{(29)} ja
\VUgras{esiliinansa} \textsuperscript{(30)}, ja hän tulee istumaan pienelle
tuolille oven eteen. Magdaleenalla on vain \VUgras{alushameensa}
\textsuperscript{(31)} ja \VUgras{korsettinsa} \textsuperscript{(32)}.

%%K t03-c1-14-1 p
14. --- Kuinka onnellinen hän onkaan! Hän unohtaa lempikukkansa,
\VUgras{neilikkansa} \textsuperscript{(34)}, joille \VUgras{perhoset}
\textsuperscript{(35)} laskeutuvat, \VUgras{tulppaaninsa}
\textsuperscript{(36)}, \VUgras{kielonsa} \textsuperscript{(37)},
\VUgras{ruukuissa} \textsuperscript{(38)}, \VUgras{muurilla}
\textsuperscript{(33)}, ja \VUgras{ruusunsa} \textsuperscript{(39)}, jotka
muodostavat niin kauniin aidan. Hän katselee saattueen naisten rikkaita
vaatteita.

%%K t03-c1-15-1 p
15. --- Kylän lapset eivät pane niin paljon merkille vaatteita. He syöksyvät
esiin ja tyrkkivät toisiaan saadakseen \VUgras{karamelleja}
\textsuperscript{(55)} tai \VUgras{rahoja} \textsuperscript{(52)}. Yksi heistä
itkee \textsuperscript{(40)}.

%%K t03-c1-16-1 p
16. --- Toinen astui \VUgras{koiran} \textsuperscript{(42)} tassulle, ja koira
pakenee kirkuen ja karkottaa \VUgras{kanan} \textsuperscript{(43)} ja sen
\VUgras{tiput} \textsuperscript{(44)}.

%%K t03-tit-6 sub
\VUsaut{5.0mm}
\VUpk{10.2pt}{40}{Kastesaattue.}
\VUsaut{4.0mm}

%%K t03-c1-17-1 p
17. --- Saattueen edellä kulkee \VUgras{imettäjä} \textsuperscript{(45)}.
Hänellä on kaunis \VUgras{myssy} \textsuperscript{(46)}, jossa on pitkät nauhat.
Hän kantaa \VUgras{vastasyntynyttä} \textsuperscript{(47)}, joka on kokonaan
puettu \VUgras{pitseihin} \textsuperscript{(48)}.

%%K t03-c1-18-1 p
18. --- Imettäjän takana tulevat \VUgras{kummisetä} \textsuperscript{(49)} ja
\VUgras{kummitäti} \textsuperscript{(53)}. He jakavat karamellit ja rahat.
Kummisedällä on \VUgras{hansikkaat} \textsuperscript{(51)} ja
\VUgras{silinterihattu} \textsuperscript{(50)}. Kummitäti pitää kaunista
\VUgras{kukkakimppua} \textsuperscript{(54)} kädessään.

%%K t03-c1-19-1 p
19. --- Heidän takanaan näemme lapsen \VUgras{sedän} \textsuperscript{(56)} ja
\VUgras{tädin} \textsuperscript{(58)}. Tädillä on tyylikäs \VUgras{hattu}
\textsuperscript{(59)}, sedällä musta \VUgras{takki} \textsuperscript{(57)} ja
valkoinen liivi. Tässä ovat sitten \VUgras{serkku} \textsuperscript{(60)} ja
\VUgras{serkku} \textsuperscript{(61)}. Tämä viimeksi mainittu pitää kädestä
vastasyntyneen \VUgras{sisarta} \textsuperscript{(62)}, pikku Bertaa.

%%K t03-c1-20-1 p
20. --- Bertan toinen \VUgras{veli} \textsuperscript{(63)} kulkee takana. Hän
antaa kätensä eräälle \VUgras{sukulaiselle} \textsuperscript{(64)}. Perheen
\VUgras{ystävät} \textsuperscript{(65)} tulevat sitten.

%%K t03-tit-7 sub
\VUsaut{4.6mm}
\VUpk{10.2pt}{40}{Katselijat.}
\VUsaut{3.4mm}

%%K t03-c1-21-1 p
21. --- Saattueen vieressä, tässä on \VUgras{kerjäläinen} \textsuperscript{(66)},
joka nojaa \VUgras{kainalosauvaansa} \textsuperscript{(67)}. Hän pyytää
almua.

%%K t03-c1-22-1 p
22. --- Toisella puolella seisoo pieni hyvin \VUgras{kohtelias}
\textsuperscript{(68)} poika. Hän tervehtii ja toivottaa
« \VUgras{hyvää päivää} » tuntemilleen henkilöille.

%%K t03-c1-23-1 p
23. --- Kyläläiset ovat tulleet ulos kodeistaan nähdäkseen kastesaattueen.
Näemme \VUgras{talonpojan} \textsuperscript{(69)}
\VUgras{puuvillamyssyineen} ja hänen vieressään
\VUgras{talonpoikaisvaimon} \textsuperscript{(70)}. Sitten \VUgras{vanhan
naisen} \textsuperscript{(71)}, joka nojaa \VUgras{keppiinsä}
\textsuperscript{(72)}. Lopuksi \VUgras{myllärin} \textsuperscript{(73)} suurine
\VUgras{huopahattuineen} \textsuperscript{(74)} ja nuoren talonpoikaisvaimon,
jolla on \VUgras{puukengät} \textsuperscript{(75)} jalassa ja joka opettaa pientä
lastaan kävelemään.

%%K t03-c1-24-1 p
24. --- Tuo kohtaus on hyvin hauska.
\VUsaut{4.0mm}

%%K t03-tit-8 sub
\VUcentre{12.0pt}{0}{\textit{Toinen kohtaus.}}
\VUsaut{3.0mm}

%%K t03-tit-9 sub
\VUpk{10.2pt}{40}{Julkinen juhla.}
\VUsaut{2.4mm}

%%K t03-tit-10 sub
\VUpk{10.2pt}{40}{Vesileikit ja soutukilpailu.}
\VUsaut{3.0mm}

%%K t03-c2-01-1 p
1. --- \VUgras{Kesä} on tullut. \VUgras{Kesäkuun} (heinäkuun tai
\VUgras{elokuun}) kuuma aurinko yllyttää nuoria kylpemään ja soutamaan.

%%K t03-c2-02-1 p
2. --- Joen oikealla rannalla soutajat riisuutuvat osallistuakseen vesileikkeihin
ja soutukilpailuun.

%%K t03-c2-03-1 p
3. --- Yksi heistä ottaa pois \VUgras{kenkänsä} \textsuperscript{(1)}; hän avaa
niiden nauhat (napit). Hänen kaksi toveriaan ovat jo valmiit. Nyt heillä on vain
\VUgras{paitansa} \textsuperscript{(2)} ja \VUgras{uimahousunsa}
\textsuperscript{(3)}. He keskustelevat odottaessaan ystäviään. He puhuvat
markkinoista ja juhlasta, jotka pidetään toisella rannalla.

%%K t03-c2-04-1 p
4. --- Maassa, heidän vieressään, ovat toverien \VUgras{vaatteet}: \VUgras{pari}
\VUgras{saapikkaita} \textsuperscript{(4)}, \VUgras{kenkiä}
\textsuperscript{(5)}, \VUgras{sukkia} \textsuperscript{(6)},
\VUgras{huopahattuja} \textsuperscript{(7)} ja \VUgras{olkihattuja}
\textsuperscript{(8)} sekä \VUgras{kaulahuivi} \textsuperscript{(9)}.

%%K t03-c2-05-1 p
5. --- Yksi nuorista on huolellisesti taittanut kokoon \VUgras{takkinsa}
\textsuperscript{(10)}. Hän panee sen pyyheliinalle, johon hänen naapurinsa pian
käärii ne kaksi \VUgras{vaatekappaletta}.

%%K t03-c2-06-1 p
6. --- Kuudes poika on myöhässä. Hänellä on yhä \VUgras{lakkinsa}
\textsuperscript{(11)} päässään. Hän ei ole ottanut pois \VUgras{paitaansa}
\textsuperscript{(12)}, eikä \VUgras{kauluksiaan} \textsuperscript{(13)}, eikä
\VUgras{kalvosimiaan} \textsuperscript{(14)}. Hän pitää \VUgras{liiviään}
\textsuperscript{(17)} kädessään, ja hänellä on yhä \VUgras{housunsa}
\textsuperscript{(16)} ja \VUgras{henkselinsä} \textsuperscript{(15)}. Hän ei
varmaankaan ehdi valmiiksi seuraavaan kilpailuun.

%%K t03-c2-07-1 p
7. --- Nuorten vieressä polku, jonka reunoilla on paljon \VUgras{ohdakkeita}
\textsuperscript{(18)}, laskeutuu joen rannalle, missä isä Jaakob
\VUgras{kaislikon} \textsuperscript{(19)} keskellä tekee valmiiksi
\VUgras{ilotulitusta} \textsuperscript{(20)}, joka on määrä laukaista illalla.

%%K t03-tit-11 sub
\VUsaut{4.4mm}
\VUpk{10.2pt}{40}{Kilpailu.}
\VUsaut{3.4mm}

%%K t03-c2-08-1 p
8. --- Joella muutamat naiset, jotka suojaavat itseään päivänvarjoillaan,
tekevät retken \VUgras{veneellä} \textsuperscript{(21)}. He seuraavat kilpailua,
joka on hyvin jännittävä. Jokaisessa \VUgras{jollassa} \textsuperscript{(22)}
soutaa neljä \VUgras{soutajaa} \textsuperscript{(24)} kaikin voimin, samalla kun
\VUgras{perämies} \textsuperscript{(26)} pitää \VUgras{peräsintä}
\textsuperscript{(25)} ja ohjaa haurasta soutuvenettä. \VUgras{Airot}
\textsuperscript{(23)} lyövät tahdissa vettä ja kevyet veneet kiitävät
huimaavaa vauhtia.

%%K t03-c2-09-1 p
9. --- Muutamat nuoret kilpailevat keskenään, sillan pilarin vieressä,
\VUgras{rasvatun maston} \textsuperscript{(28)} palkinnosta. Yksi heistä kävelee
varovasti taipuisaa ja liukasta mastoa pitkin päästäkseen pieneen
\VUgras{lippuun} \textsuperscript{(29)}, joka on kiinnitetty sen päähän. Hänen
onneton \VUgras{kilpakumppaninsa} \textsuperscript{(30)} on pudonnut veteen. Hän
ui päästäkseen maihin.

%%K t03-c2-10-1 p
10. --- Vanhemmat nuoret pitävät \VUgras{(vene)turnajaisia}
\textsuperscript{(32)} \VUgras{laiturin} \textsuperscript{(33)} luona. Kaikkia
näitä leikkejä seuraa lukuisa \VUgras{yleisö} \textsuperscript{(31)}, joka nojaa
\VUgras{sillan} \textsuperscript{(27)} \VUgras{kaiteeseen} ja tungeksii
\VUgras{rannalla}. Muutamat katselijat kääntyvät kuitenkin ympäri seuratakseen
leikkiä \VUgras{palkintomastolla} \textsuperscript{(45)} tai ihaillakseen
\VUgras{kaupunginhallituksen saattuetta}, joka tulee
\VUgras{palkintojenjakoon}.

%%K t03-c2-11-1 p
11. --- Ensin, tässä on muutamia \VUgras{katupoikia} \textsuperscript{(35)},
jotka kulkevat \VUgras{soittajien} \textsuperscript{(36)} edellä; sitten tulee
\VUgras{pormestari} \textsuperscript{(37)}, hänen sijaisensa ja pienen kaupungin
\VUgras{kaupunginvaltuutetut}. Lopuksi kulkevat \VUgras{palomiehet}
\textsuperscript{(38)}.

%%K t03-tit-12 sub
\VUsaut{5.0mm}
\VUpk{10.2pt}{40}{Markkinat.}
\VUsaut{3.6mm}

%%K t03-c2-12-1 p
12. --- Lukuisia ihmisiä on \VUgras{markkinapaikalla} \textsuperscript{(39)}.
Tullaan katsomaan eri \VUgras{kojuja}: \VUgras{puuhevosia}
\textsuperscript{(40)}, \VUgras{sirkusta} \textsuperscript{(41)}, jossa näytäntö
on jo alkanut; \VUgras{julkisia tanssiaisia} \textsuperscript{(42)}, joissa
kaupungin nuoret miehet ja nuoret naiset nauttivat \VUgras{tanssin} ilosta.

%%K t03-c2-13-1 p
13. --- Perällä näemme \VUgras{areenan} \textsuperscript{(44)}, jossa urhea
espanjalainen \VUgras{matadori} taistelee raivostunutta \VUgras{härkää}
\textsuperscript{(45)} vastaan, sitten \VUgras{rullaluistinradan}
\textsuperscript{(49)}, \VUgras{ampumaradan} \textsuperscript{(46)}, jossa
harjoitellaan \VUgras{kiväärien} käyttöä ja ampumista \VUgras{maalitauluun}.

%%K t03-c2-14-1 p
14. --- Aivan lähellä, tässä on \VUgras{markkinateatteri} \textsuperscript{(47)},
jossa esitetään \VUgras{komediaa} ja \VUgras{draamaa}. Aivan vieressä ovat
\VUgras{jättiläisen} \textsuperscript{(48)} ja \VUgras{kääpiön} kojut.
Ensimmäisen \VUgras{pituus} on kaksi metriä ja viisitoista senttimetriä; toinen
on tuskin niin suuri kuin lapsi.

%%K t03-c2-15-1 p
15. --- Lopuksi, tässä on suuri \VUgras{eläintarha} \textsuperscript{(50)}
\VUgras{villipetoineen}, \VUgras{tiikereineen} \textsuperscript{(51)}, joilla on
voimakkaat \VUgras{kynnet}, \VUgras{leijonineen} \textsuperscript{(52)}, jolla
on tuuhea \VUgras{harja}, kahtine \VUgras{kirahveineen} \textsuperscript{(53)}
ja \VUgras{norsuineen} \textsuperscript{(54)}, joka niin taitavasti käyttää
\VUgras{kärsäänsä}.

%%K t03-c2-16-1 p
16. --- \VUgras{Eläintenkesyttäjä} \textsuperscript{(55)}, joka kouluttaa noita
eläimiä, seisoo kojun ovella.
\VUsaut{6.0mm}
\VUfilet{22mm}
